\chapter{Families with twelve linear factors}\label{sec:families}

In this section, we set up the framework shared by the three families of \cref{sec:rank5,sec:rank6,sec:rank7} and prove once the steps that are common to them. A twelve-form datum (\cref{def:datum}) consists of twelve binary linear forms and three integers $A$, $B$, $C$ satisfying the identity $A^2F+4BH=CJ$, where $F$, $H$, $J$ are products of four of the forms each. For each $(M,N)\in\Z^2$ with $F(M,N)H(M,N)J(M,N)\ne0$, it gives a pair of $2$-isogenous curves $E_{M,N}$ and $E'_{M,N}$. In \cref{sec:families-data}, we compute their basic invariants (\cref{lem:datum-basic,lem:j-invariant}). In \cref{sec:families-prime}, we introduce prime specializations, at which the twelve forms take prime values up to fixed cofactors (\cref{def:prime-specialization}), show that a congruence condition on $(M,N)$ forces the divisibility $c_i\mid L_i(M,N)$ and fixes the residues of the quotients (\cref{lem:precision}), and determine the local images of the descent maps at these twelve primes (\cref{lem:moving-images}). \Cref{sec:families-rows} turns these local images into linear conditions (\cref{lem:moving-rows}). The coefficients of these conditions are quadratic characters, and \cref{sec:families-reciprocity} shows that quadratic reciprocity determines them through fixed residues and a single sign $\eps$ (\cref{lem:reciprocity,cor:eta-structure}). Finally, \cref{sec:families-criterion} proves the rank criterion (\cref{prop:rank-criterion}), which reduces the rank of $E_{M,N}(\Q)$ at all prime specializations in a given set to finitely many checks, for both values of $\eps$.

This section uses the results of \cref{sec:descent}. \Cref{sec:primes} uses \cref{def:datum,def:prime-specialization,lem:prime-spec-basic,lem:j-invariant} to produce infinitely many prime specializations, and \cref{sec:rank5,sec:rank6,sec:rank7} verify the hypotheses of \cref{prop:rank-criterion} for the three families.

\section{Twelve-form data}\label{sec:families-data}

\begin{definition}[Twelve-form datum]\label{def:datum}
A \emph{twelve-form datum} $\datum$ consists of the following.
\begin{enumerate}[label=(\alph*),ref=(\alph*)]
\item\label{item:datum-forms} Pairs $(a_i,b_i)\in\Z^2$ for $0\le i\le11$, with $\gcd(a_i,b_i)=1$ for every $i$, $(a_0,b_0)=(1,0)$, $b_i\ne0$ for $1\le i\le11$, and $\Delta_{ij}=a_ib_j-a_jb_i\ne0$ for all $i\ne j$.
\item\label{item:datum-identity} Nonzero integers $A$, $B$, $C$ such that
\begin{equation}\label{eq:datum-identity}
A^2F+4BH=CJ
\end{equation}
in $\Z[M,N]$, where $L_i=a_iM+b_iN$ for $0\le i\le11$ and
\begin{equation}\label{eq:FHJ}
F=L_0L_1L_2L_3,\qquad H=L_4L_5L_6L_7,\qquad J=L_8L_9L_{10}L_{11}.
\end{equation}
\item\label{item:datum-cofactors} Positive integers $c_0,\dots,c_{11}$, the \emph{cofactors}, with $\min_ic_i=1$.
\item\label{item:datum-constants} Two finite ordered lists $\Gamma_E$ and $\Gamma_{E'}$ of nonzero integers.
\end{enumerate}
For $(M,N)\in\Z^2$ with $F(M,N)H(M,N)J(M,N)\ne0$, the \emph{curves} of $\datum$ are
\begin{equation}\label{eq:family-curves}
E_{M,N}\colon y^2=x^3+AFx^2-BFHx,\qquad E'_{M,N}\colon y^2=x^3-2AFx^2+CFJx,
\end{equation}
where $F$, $H$, $J$ are evaluated at $(M,N)$.
\end{definition}

Throughout this section, we fix a twelve-form datum $\datum$. We write
\begin{equation}\label{eq:index-sets}
I_F=\{0,1,2,3\},\qquad I_H=\{4,5,6,7\},\qquad I_J=\{8,9,10,11\},
\end{equation}
so that $F=\prod_{i\in I_F}L_i$, $H=\prod_{i\in I_H}L_i$ and $J=\prod_{i\in I_J}L_i$. The moving index sets and the opposite index sets of the two curves are
\begin{equation}\label{eq:moving-sets}
G_E=I_F\cup I_H,\qquad G_{E'}=I_F\cup I_J,\qquad O_E=I_J,\qquad O_{E'}=I_H.
\end{equation}
The symbol $\star$ stands for either $E$ or $E'$. When $(M,N)$ is fixed, we write $E=E_{M,N}$ and $E'=E'_{M,N}$, and we write
\begin{equation}\label{eq:family-coefficients}
\alpha=AF,\qquad\beta=-BFH,\qquad\alpha'=-2AF,\qquad\beta'=CFJ
\end{equation}
for their coefficients, so that $E=E_{\alpha,\beta}$ and $E'=E_{\alpha',\beta'}$ in the notation of \eqref{eq:E-alpha-beta}. We let $\delta_E$, $\delta_{E'}$ be their descent maps (\cref{def:descent-map}), $\delta_{E,v}$, $\delta_{E',v}$ the local descent maps, and
\begin{equation}\label{eq:dE-dEprime}
d_E=\dim_{\F_2}\delta_E(E(\Q)),\qquad d_{E'}=\dim_{\F_2}\delta_{E'}(E'(\Q)).
\end{equation}
Condition \ref{item:datum-forms} of \cref{def:datum} says that the twelve forms are primitive and pairwise nonproportional, so that their zeros are twelve distinct points of $\mathbb P^1$. The cofactors and the lists $\Gamma_E$, $\Gamma_{E'}$ are used only from \cref{sec:families-prime} on.

\begin{lemma}[Basic properties of the curves]\label{lem:datum-basic}
Let $(M,N)\in\Z^2$ with $FHJ\ne0$.
\begin{enumerate}[label=(\alph*),ref=(\alph*)]
\item\label{item:datum-basic-dual} We have $\alpha^2-4\beta=CFJ=\beta'$ and $\alpha'^2-4\beta'=16\beta$. In particular, $E'_{M,N}$ is the curve $E_{\alpha',\beta'}$ of \eqref{eq:isogenous-coefficients}, and the maps $\phi$ and $\dual$ of \eqref{eq:isogenies} are the $2$-isogenies between $E_{M,N}$ and $E'_{M,N}$.
\item\label{item:datum-basic-disc} The discriminants of the models \eqref{eq:family-curves} are
\begin{equation}\label{eq:family-discriminants}
\Delta(E_{M,N})=16B^2CF^3H^2J,\qquad\Delta(E'_{M,N})=-256BC^2F^3HJ^2.
\end{equation}
Both are nonzero, and $E_{M,N}$ and $E'_{M,N}$ are elliptic curves over $\Q$.
\item\label{item:datum-basic-scaling} Let $(M_1,N_1)=(\lambda M,\lambda N)\in\Z^2$ with $\lambda\in\Q^\times$. Then $(x,y)\mapsto(\lambda^4x,\lambda^6y)$ is an isomorphism over $\Q$ from $E_{M,N}$ to $E_{M_1,N_1}$, and from $E'_{M,N}$ to $E'_{M_1,N_1}$. It preserves the class $\sq x$ of every point with $x\ne0$.
\end{enumerate}
\end{lemma}

\begin{proof}
By \eqref{eq:datum-identity}, we have
\begin{equation}
\alpha^2-4\beta=A^2F^2+4BFH=F(A^2F+4BH)=CFJ,
\end{equation}
and then $\alpha'^2-4\beta'=4\alpha^2-4(\alpha^2-4\beta)=16\beta$. This proves \ref{item:datum-basic-dual}. For a curve $y^2=x^3+ax^2+bx$, the standard formulas \cite[Section~III.1]{SilvermanAEC} give
\begin{equation}\label{eq:c4-Delta}
c_4=16(a^2-3b),\qquad\Delta=16b^2(a^2-4b).
\end{equation}
For $E_{M,N}$, we obtain $\Delta=16\beta^2(\alpha^2-4\beta)=16B^2F^2H^2\cdot CFJ$, and for $E'_{M,N}$, we obtain $\Delta=16\beta'^2(\alpha'^2-4\beta')=16C^2F^2J^2\cdot16\beta=-256BC^2F^3HJ^2$. These are nonzero because $B$, $C$ and $FHJ$ are nonzero. For \ref{item:datum-basic-scaling}, since each of $F$, $H$, $J$ is homogeneous of degree four, the coefficients of $E_{M_1,N_1}$ are $\lambda^4\alpha$ and $\lambda^8\beta$. Then
\begin{equation}
(\lambda^6y)^2-(\lambda^4x)^3-\lambda^4\alpha(\lambda^4x)^2-\lambda^8\beta(\lambda^4x)=\lambda^{12}(y^2-x^3-\alpha x^2-\beta x),
\end{equation}
and the same computation applies to $E'$. Finally, $\sq{\lambda^4x}=\sq x$.
\end{proof}

\begin{lemma}[The $j$-invariant]\label{lem:j-invariant}
Let $\mathcal N$ and $\mathcal D$ be the binary forms of degree $12$ given by
\begin{equation}\label{eq:j-forms}
\mathcal N=256(A^2F+3BH)^3,\qquad\mathcal D=B^2CH^2J.
\end{equation}
\begin{enumerate}[label=(\alph*),ref=(\alph*)]
\item\label{item:j-formula} For $(M,N)\in\Z^2$ with $FHJ\ne0$, we have $j(E_{M,N})=\mathcal N(M,N)/\mathcal D(M,N)$.
\item\label{item:j-degree} The forms $\mathcal N$ and $\mathcal D$ have no common zero in $\mathbb P^1(\Qbar)$, and $[M:N]\mapsto[\mathcal N(M,N):\mathcal D(M,N)]$ is a morphism $\mathbb P^1\to\mathbb P^1$ of degree $12$. For every $j_0\in\Q$, at most $12$ points $[M:N]\in\mathbb P^1(\Q)$ with $FHJ(M,N)\ne0$ satisfy $j(E_{M,N})=j_0$.
\item\label{item:j-infinite} Let $\mathcal Q\subseteq\Z^2$ be a set of pairs with $FHJ\ne0$ at every point of $\mathcal Q$, such that the points $[M:N]$ with $(M,N)\in\mathcal Q$ are infinitely many. Then the curves $E_{M,N}$ with $(M,N)\in\mathcal Q$ lie in infinitely many isomorphism classes over $\Qbar$.
\end{enumerate}
\end{lemma}

\begin{proof}
By \eqref{eq:c4-Delta}, we have $j=c_4^3/\Delta=256(a^2-3b)^3/(b^2(a^2-4b))$ for the curve $y^2=x^3+ax^2+bx$. For $E_{M,N}$, we have $\alpha^2-3\beta=A^2F^2+3BFH=F(A^2F+3BH)$, $\beta^2=B^2F^2H^2$ and $\alpha^2-4\beta=CFJ$. Inserting these into the formula for $j$, we obtain
\begin{equation}
j(E_{M,N})=\frac{256F^3(A^2F+3BH)^3}{B^2F^2H^2\cdot CFJ}=\frac{256(A^2F+3BH)^3}{B^2CH^2J}.
\end{equation}
This proves \ref{item:j-formula}.

The zero of $L_i$ in $\mathbb P^1(\Qbar)$ is the point $[b_i:-a_i]$, and these twelve points are distinct because $\Delta_{ij}\ne0$ for $i\ne j$. Since $B,C\ne0$, the zeros of $\mathcal D$ are the zeros of $L_i$ with $i\in I_H\cup I_J$. At a zero of $L_i$ with $i\in I_H$, we have $H=0$ and $F\ne0$, which gives $A^2F+3BH=A^2F\ne0$. At a zero of $L_i$ with $i\in I_J$, we have $J=0$ and $H\ne0$, and \eqref{eq:datum-identity} implies that $A^2F=-4BH$; then $A^2F+3BH=-BH\ne0$. We conclude that $\mathcal N$ and $\mathcal D$ have no common zero. As both have degree $12$, the map has degree $12$. Let $j_0\in\Q$. The form $\mathcal N-j_0\mathcal D$ is not zero, since otherwise $\mathcal N$ would vanish at the zeros of $\mathcal D$. As a nonzero binary form of degree $12$, it has at most $12$ zeros in $\mathbb P^1(\Qbar)$. If $FHJ(M,N)\ne0$ and $j(E_{M,N})=j_0$, then $\mathcal D(M,N)\ne0$, and \ref{item:j-formula} implies that $(\mathcal N-j_0\mathcal D)(M,N)=0$. This proves \ref{item:j-degree}.

For \ref{item:j-infinite}, two elliptic curves over $\Qbar$ are isomorphic if and only if they have the same $j$-invariant \cite[Proposition~III.1.4(b)]{SilvermanAEC}. If the curves $E_{M,N}$ with $(M,N)\in\mathcal Q$ lay in finitely many classes, then $j(E_{M,N})$ would take finitely many values on $\mathcal Q$, and by \ref{item:j-degree} the points $[M:N]$ would be finitely many.
\end{proof}

\begin{remark}[The $j$-invariants of the three families]\label{rem:j-families}
For the data of \cref{sec:rank5,sec:rank6,sec:rank7}, \cref{lem:j-invariant}\ref{item:j-formula} gives the following $j$-invariants; in the rank-seven family, $33124=182^2$.
\begin{center}
\begin{tabular}{lcc}
\toprule
Family & $(A,B,C)$ & $j(E_{M,N})$\\
\midrule
Rank five & $(1500,2700,57600)$ & $\dfrac49\cdot\dfrac{(2500F+9H)^3}{H^2J}$\\[8pt]
Rank six & $(70,784,1764)$ & $\dfrac{16}9\cdot\dfrac{(25F+12H)^3}{H^2J}$\\[8pt]
Rank seven & $(91205,33124,9)$ & $\dfrac{256}{9\cdot182^4}\cdot\dfrac{(91205^2F+3\cdot182^2H)^3}{H^2J}$\\
\bottomrule
\end{tabular}
\end{center}
For example, in the rank-five family $A^2F+3BH=900(2500F+9H)$ and $256\cdot900^3/(2700^2\cdot57600)=4/9$.
\end{remark}

\begin{remark}[The associated elliptic surfaces]\label{rem:elliptic-surfaces}
The homogeneous equations define Weierstrass models over the projective parameter line, with coefficients $\alpha$ and $\beta$ of degrees four and eight. In either projective chart, the formulas in \eqref{eq:c4-Delta} give
\[
c_4(E)=16F(A^2F+3BH),\qquad \Delta(E)=16B^2CF^3H^2J.
\]
The twelve forms have distinct zeros. At the zeros of $F$, $H$ and $J$, the orders of $(c_4,\Delta)$ are $(1,3)$, $(0,2)$ and $(0,1)$, respectively; at a zero of $J$, we use $A^2F+3BH=-BH$. These models are minimal because the discriminant orders are less than $12$ \cite[p.~40, Remark~1]{Tate1975}. Tate's algorithm gives four fibers of type $\mathrm{III}$, four of type $\mathrm I_2$ and four of type $\mathrm I_1$ \cite[p.~34]{Tate1975}; see also \cite[\S4.5, Table~1]{SchuttShioda2010}. For $E'$, the types at $H$ and $J$ are interchanged. There are no other singular fibers, and the minimal discriminant divisor has degree $24$. Thus the relatively minimal smooth elliptic surfaces have $\chi(\mathcal O)=2$. Over the base $\mathbb P^1$, the canonical bundle formula and $h^1(\mathcal O)=0$ then show that they are K3 surfaces \cite[Theorems~6.8 and~6.12 and \S8.2]{SchuttShioda2010}.
\end{remark}

\section{Prime specializations}\label{sec:families-prime}

\begin{definition}[Prime specializations]\label{def:prime-specialization}
Let $S$ be a finite set of primes that contains every prime dividing
\begin{equation}\label{eq:S-condition}
6ABC\prod_{i=0}^{11}c_i\prod_{i=1}^{11}b_i\prod_{0\le i<j\le11}\Delta_{ij}
\end{equation}
and every prime dividing an entry of $\Gamma_E$ or $\Gamma_{E'}$. A pair $(M,N)\in\Z^2$ is a \emph{prime specialization} of $\datum$ (with respect to $S$) if $c_i$ divides $L_i(M,N)$ for $0\le i\le11$ and the twelve numbers
\begin{equation}\label{eq:prime-values}
p_i=\frac{L_i(M,N)}{c_i},\qquad0\le i\le11,
\end{equation}
are distinct primes, none of which lies in $S$.
\end{definition}

In \cref{sec:rank5,sec:rank6,sec:rank7}, the divisibility $c_i\mid L_i(M,N)$ and the residues of the numbers $L_i(M,N)/c_i$ come from a congruence condition on $(M,N)$, through the following elementary fact.

\begin{lemma}[Divisibility and residues on a congruence class]\label{lem:precision}
Let $L=aM+bN$ with $a,b\in\Z$, let $W_0$ be a positive integer, and let $(m,n)\in\Z^2$. Let $c$ be a positive integer that divides both $W_0$ and $L(m,n)$. Then every $(M,N)\in\Z^2$ with $(M,N)\equiv(m,n)\pmod{W_0}$ satisfies $c\mid L(M,N)$ and
\begin{equation}\label{eq:precision}
\frac{L(M,N)}c\equiv\frac{L(m,n)}c\pmod{W_0/c}.
\end{equation}
In particular, if $q$ is a prime that divides $W_0/c$ and does not divide $L(m,n)/c$, then $v_q(L(M,N))=v_q(c)$.
\end{lemma}

\begin{proof}
Since $M-m$ and $N-n$ are divisible by $W_0$, there exists $s\in\Z$ such that $L(M,N)=L(m,n)+W_0s$. Dividing by $c$, we obtain
\begin{equation}
\frac{L(M,N)}c=\frac{L(m,n)}c+\frac{W_0}c\,s,
\end{equation}
where both terms on the right are integers. This proves that $c$ divides $L(M,N)$ and that \eqref{eq:precision} holds. If $q\mid W_0/c$ and $q\nmid L(m,n)/c$, then \eqref{eq:precision} shows that $L(M,N)/c$ is an integer prime to $q$, which implies that $v_q(L(M,N))=v_q(c)$.
\end{proof}

At a prime specialization, we write $\lambda_i$ for the restriction of the character $\chi_{p_i}$ of \eqref{eq:chi-def} to the integers prime to $p_i$, so that $\leg z{p_i}=(-1)^{\lambda_i(z)}$ for such integers $z$. We set
\begin{equation}\label{eq:kappa}
\kappa_E=-B\prod_{j\in G_E}c_j,\qquad\kappa_{E'}=C\prod_{j\in G_{E'}}c_j.
\end{equation}
The \emph{ambient list} $\cB_\star$ consists of the entries of $\Gamma_\star$, in their order, followed by the primes $p_j$ with $j\in G_\star$, in increasing order of $j$; its length is $n_\star=|\Gamma_\star|+8$. We index the entries of a vector in $\F_2^{n_\star}$ by the entries $g$ of $\Gamma_\star$, which we call the constant columns, and by the indices $j\in G_\star$, which we call the moving columns. We let $u_\star\in\F_2^{n_\star}$ be the vector with entry $0$ at each constant column and $1$ at each moving column.

\begin{lemma}[Basic properties of prime specializations]\label{lem:prime-spec-basic}
Let $(M,N)$ be a prime specialization.
\begin{enumerate}[label=(\alph*),ref=(\alph*)]
\item\label{item:psb-nonzero} The product $FHJ$ is nonzero, and $E=E_{M,N}$ and $E'=E'_{M,N}$ are elliptic curves. Each $p_i$ is odd, and it divides none of the integers $A$, $B$, $C$, $c_j$ (for all $j$), $b_j$ (for $j\ge1$) and $\Delta_{jk}$ (for $j\ne k$), and no entry of $\Gamma_E$ or $\Gamma_{E'}$.
\item\label{item:psb-valuations} For all $i$ and $j$, the valuation $v_{p_i}(L_j)$ is $1$ if $i=j$ and $0$ if $i\ne j$.
\item\label{item:psb-gcd} The integers $M$ and $N$ are coprime, $M=c_0p_0$, and $p_0\nmid N$.
\item\label{item:psb-torsion} The valuation $v_{p_0}(CFJ)$ is $1$. In particular, $CFJ=\alpha^2-4\beta$ is not a square in $\Q$, and $E(\Q)[2]=\{O,T\}$.
\item\label{item:psb-beta} The coefficients satisfy $\beta=\kappa_E\prod_{j\in G_E}p_j$ and $\beta'=\kappa_{E'}\prod_{j\in G_{E'}}p_j$.
\item\label{item:psb-independent} If $\Gamma_\star$ consists of $-1$ and distinct primes, then the classes of the entries of $\cB_\star$ are linearly independent in $\sqgp\Q$.
\item\label{item:psb-support} If $\Gamma_\star$ contains $-1$ and every prime dividing $\kappa_\star$, then $\delta_\star(\star(\Q))$ is contained in the span of the classes of the entries of $\cB_\star$.
\end{enumerate}
In particular, if $\Gamma_\star$ consists of $-1$ and distinct primes, including every prime dividing $\kappa_\star$, then $\cB_\star$ is an ambient basis for $\star$ in the sense of \cref{def:rows}.
\end{lemma}

\begin{proof}
We have $L_i(M,N)=c_ip_i\ne0$ for every $i$. Then $FHJ\ne0$, and \cref{lem:datum-basic}\ref{item:datum-basic-disc} shows that $E$ and $E'$ are elliptic curves. The prime $2$ divides \eqref{eq:S-condition} and lies in $S$, and each $p_i$ is odd. The rest of \ref{item:psb-nonzero} holds because $p_i\notin S$. For \ref{item:psb-valuations}, we have $L_j=c_jp_j$ with $p_{i}\nmid c_j$ by \ref{item:psb-nonzero}, and $p_i\ne p_j$ for $j\ne i$.

For \ref{item:psb-gcd}, let $q$ be a prime dividing both $M$ and $N$. Then $q$ divides $L_i=c_ip_i$ for every $i$. By \cref{def:datum}\ref{item:datum-cofactors}, there exists $i_1$ with $c_{i_1}=1$, and then $q=p_{i_1}$. For $i\ne i_1$, we have $q\ne p_i$, so that $q$ divides $c_i$; this contradicts \ref{item:psb-nonzero} for the prime $p_{i_1}$. Hence $\gcd(M,N)=1$. Since $(a_0,b_0)=(1,0)$, we have $M=L_0=c_0p_0$, and then $p_0\nmid N$.

For \ref{item:psb-torsion}, part \ref{item:psb-valuations} implies that $v_{p_0}(F)=1$ and $v_{p_0}(J)=0$, and $p_0\nmid C$ by \ref{item:psb-nonzero}. Then $v_{p_0}(CFJ)=1$ is odd, and $CFJ$ is not a square in $\Q$. By \cref{lem:datum-basic}\ref{item:datum-basic-dual}, $CFJ=\alpha^2-4\beta$, and \cref{lem:two-torsion} with $K=\Q$ shows that $T$ is the only point of order two in $E(\Q)$.

Part \ref{item:psb-beta} follows from $L_j=c_jp_j$ and the definitions \eqref{eq:family-coefficients} and \eqref{eq:kappa}. For \ref{item:psb-independent}, suppose that a product of entries of $\cB_\star$ is a square in $\Q$. For $j\in G_\star$, the prime $p_j$ divides no other entry of $\cB_\star$, by \ref{item:psb-nonzero} and since the $p_i$ are distinct. The valuation at $p_j$ then shows that $p_j$ occurs to an even power. The valuations at the primes of $\Gamma_\star$ show the same for these primes, and then the sign shows it for $-1$.

For \ref{item:psb-support}, the coefficients $\alpha,\beta$ and $\alpha',\beta'$ are integers with $\beta(\alpha^2-4\beta)\ne0$ and $\beta'(\alpha'^2-4\beta')\ne0$. By \cref{lem:support}, every class in $\delta_E(E(\Q))$ is the class of $\pm1$ times a product of distinct primes dividing $\beta$, and by \ref{item:psb-beta} these primes divide $\kappa_E$ or are equal to $p_j$ with $j\in G_E$. The same argument applies to $E'$ with $\beta'$ and $\kappa_{E'}$. The final statement combines \ref{item:psb-independent} and \ref{item:psb-support}.
\end{proof}

\begin{lemma}[Local images at the moving primes]\label{lem:moving-images}
Let $(M,N)$ be a prime specialization, let $0\le i\le11$, and let $p=p_i$.
\begin{enumerate}[label=(\alph*),ref=(\alph*)]
\item\label{item:moving-F} If $i\in I_F$, then $\delta_{E,p}(E(\Q_p))=\{\sq1,\sq\beta\}$ and $\delta_{E',p}(E'(\Q_p))=\{\sq1,\sq{\beta'}\}$.
\item\label{item:moving-J} If $i\in I_J$, then $\delta_{E,p}(E(\Q_p))=\{\sq1\}$.
\item\label{item:moving-H} If $i\in I_H$, then $\delta_{E',p}(E'(\Q_p))=\{\sq1\}$.
\end{enumerate}
\end{lemma}

\begin{proof}
By \cref{lem:prime-spec-basic}\ref{item:psb-nonzero} and \ref{item:psb-valuations}, the prime $p$ is odd and does not divide $ABC$. Further, $v_p(F)=1$ if $i\in I_F$ and $v_p(F)=0$ otherwise, and similarly for $H$ with $I_H$ and for $J$ with $I_J$. By \eqref{eq:family-coefficients} and \cref{lem:datum-basic}\ref{item:datum-basic-dual}, the valuations of the coefficients are as follows.
\begin{center}
\begin{tabular}{lcccl}
\toprule
$i\in$ & $v_p(\alpha)=v_p(\alpha')$ & $v_p(\beta)$ & $v_p(\beta')$ & Case of \cref{lem:odd-local}\\
\midrule
$I_F$ & $1$ & $1$ & $1$ & \ref{item:odd-local-F} for $E$ and for $E'$\\
$I_J$ & $0$ & $0$ & $1$ & \ref{item:odd-local-J} for $E$, since $\alpha^2-4\beta=\beta'$\\
$I_H$ & $0$ & $1$ & $0$ & \ref{item:odd-local-J} for $E'$, since $\alpha'^2-4\beta'=16\beta$\\
\bottomrule
\end{tabular}
\end{center}
If $i\in I_F$, then $v_p(\alpha)\ge1$ and $v_p(\beta)=1$, and \cref{item:odd-local-F} shows that $\delta_{E,p}(E(\Q_p))=\{\sq1,\sq\beta\}$; the same holds for $E'$, with $\alpha',\beta'$ in place of $\alpha,\beta$. If $i\in I_J$, then $\alpha$ and $\beta$ are units and $v_p(\alpha^2-4\beta)=v_p(\beta')=1$, and \cref{item:odd-local-J} gives \ref{item:moving-J}. If $i\in I_H$, then $\alpha'$ and $\beta'$ are units and $v_p(\alpha'^2-4\beta')=v_p(16\beta)=1$, and \ref{item:moving-H} follows from \cref{item:odd-local-J}.
\end{proof}

\section{Rows at the moving primes}\label{sec:families-rows}

The following rows encode \cref{lem:moving-images} in the coordinates of \cref{lem:squares}\ref{item:squares-odd}. By \cref{lem:prime-spec-basic}\ref{item:psb-nonzero} and \ref{item:psb-valuations}, every integer to which we apply $\lambda_i$ below is prime to $p_i$.

\begin{definition}[Moving rows]\label{def:moving-rows}
Let $(M,N)$ be a prime specialization and $\star\in\{E,E'\}$. For $i\in I_F$, the \emph{moving row} $r^\star_i\in\F_2^{n_\star}$ has the entries
\begin{equation}\label{eq:moving-row-F}
\begin{array}{ll}
\lambda_i(g) & \text{at each constant column }g,\\
\lambda_i(p_j) & \text{at each moving column }j\in G_\star\setminus\{i\},\\
\lambda_i(\kappa_\star)+\sum_{j\in G_\star\setminus\{i\}}\lambda_i(p_j)\quad & \text{at the moving column }i.
\end{array}
\end{equation}
For $i\in O_\star$, the \emph{moving row} $r^\star_i$ has entry $\lambda_i(g)$ at each constant column $g$ and entry $\lambda_i(p_j)$ at each moving column $j\in G_\star$.
\end{definition}

\begin{lemma}[Conditions at the moving primes]\label{lem:moving-rows}
Let $(M,N)$ be a prime specialization, $\star\in\{E,E'\}$ and $i\in I_F\cup O_\star$. Then $r^\star_i$ is a row for $\star$ and $\cB_\star$ at the place $p_i$, in the sense of \cref{def:rows}. In particular, $r^\star_i\cdot\mathbf w=0$ for every $\mathbf w\in\F_2^{n_\star}$ such that $\sq{\mathbf w}_{\cB_\star}\in\delta_\star(\star(\Q))$.
\end{lemma}

\begin{proof}
Let $p=p_i$. We use the coordinates $\sq z\mapsto(\nu_p(z),\chi_p(z))$ of \eqref{eq:odd-coordinates}, so that $\sq{p^ku_0}\mapsto(k\bmod2,\chi_p(u_0))$ for $u_0\in\Z_p^\times$. Every entry of $\cB_\star$ other than $p_i$ is a $p$-adic unit, and its class has coordinates $(0,\lambda_i(g))$ or $(0,\lambda_i(p_j))$; if $p_i$ is an entry of $\cB_\star$, its class has coordinates $(1,0)$.

Let $i\in I_F$, and take $\star=E$; for $\star=E'$, replace $\beta$ and $\kappa_E$ by $\beta'$ and $\kappa_{E'}$. By \cref{lem:moving-images}\ref{item:moving-F}, the local image $\delta_{E,p}(E(\Q_p))$ is the subgroup $U_p=\{\sq1,\sq\beta\}$. By \cref{lem:prime-spec-basic}\ref{item:psb-beta}, we have $\beta=p_iw$ with $w=\kappa_E\prod_{j\in G_E\setminus\{i\}}p_j$ a $p$-adic unit. The class $\sq\beta$ has coordinates $(1,h)$ with
\begin{equation}
h=\lambda_i(w)=\lambda_i(\kappa_E)+\sum_{j\in G_E\setminus\{i\}}\lambda_i(p_j).
\end{equation}
The linear map $f(k,\ell)=\ell+hk$ vanishes on $U_p=\{(0,0),(1,h)\}$. Its values at the classes of the entries of $\cB_E$ are $\lambda_i(g)$ at a constant column $g$, $\lambda_i(p_j)$ at a moving column $j\ne i$, and $h$ at the moving column $i$. These are the entries \eqref{eq:moving-row-F}, and $r^E_i$ is the row defined by $(p_i,U_p,f)$.

Let $i\in O_\star$. Then $i\notin G_\star$, and \cref{lem:moving-images}\ref{item:moving-J} (for $\star=E$) or \ref{item:moving-H} (for $\star=E'$) shows that the local image at $p$ is $U_p=\{\sq1\}$. The map $f(k,\ell)=\ell$ vanishes on $U_p$, and its values at the entries of $\cB_\star$ are the entries of $r^\star_i$. (The other coordinate $k$ vanishes at every entry of $\cB_\star$ and gives the zero row.) The last statement of the lemma is \cref{lem:upper-bound}.
\end{proof}

\section{Reciprocity and the sign \texorpdfstring{$\eps$}{epsilon}}\label{sec:families-reciprocity}

The entries of the moving rows are characters $\lambda_i(g)$ and $\lambda_i(\kappa_\star)$ of fixed integers, and characters $\lambda_i(p_j)$ of one moving prime at another. We now express the latter through characters of fixed integers and the single sign
\begin{equation}\label{eq:eps-def}
\eps=\leg N{p_0}=(-1)^\eta,\qquad\eta\in\F_2,
\end{equation}
which is defined at every prime specialization because $p_0\nmid N$ (\cref{lem:prime-spec-basic}\ref{item:psb-gcd}). We use quadratic reciprocity and its two supplementary laws \cite[Chapter~5, Section~2, Theorem~1, p.~53]{IrelandRosen}. In the notation \eqref{eq:chi-def}, they state that, for distinct odd primes $p$ and $q$,
\begin{equation}\label{eq:reciprocity}
\begin{gathered}
\chi_p(-1)=\frac{p-1}2\bmod2,\qquad\chi_p(2)=\frac{p^2-1}8\bmod2,\\
\chi_p(q)+\chi_q(p)=\chi_p(-1)\chi_q(-1).
\end{gathered}
\end{equation}

The determinant elimination and one-sign reduction below are analogous to the Gaussian reciprocity calculation in \cite[Lemma~4.1]{Savoie2026}.

\begin{lemma}[Reciprocity for the prime values]\label{lem:reciprocity}
Let $(M,N)$ be a prime specialization, and let $\eps=\leg N{p_0}$.
\begin{enumerate}[label=(\alph*),ref=(\alph*)]
\item\label{item:rec-pivot} For $1\le i\le11$, the symbol $\leg{p_i}{p_0}$ is equal to $\leg{b_ic_i}{p_0}\eps$.
\item\label{item:rec-reverse} For $1\le i\le11$, we have $\leg{p_0}{p_i}=(-1)^{(p_0-1)(p_i-1)/4}\leg{p_i}{p_0}$.
\item\label{item:rec-cross} For distinct $i,j$ with $1\le i,j\le11$, we have $b_iL_j\equiv\Delta_{ji}M\pmod{p_i}$, and hence
\begin{equation}\label{eq:rec-cross}
\leg{p_j}{p_i}=\leg{c_0c_jb_i\Delta_{ji}}{p_i}\leg{p_0}{p_i}.
\end{equation}
\end{enumerate}
\end{lemma}

\begin{proof}
For \ref{item:rec-pivot}, recall that $M=c_0p_0$ (\cref{lem:prime-spec-basic}\ref{item:psb-gcd}). Reducing $c_ip_i=a_iM+b_iN$ modulo $p_0$, we find $c_ip_i\equiv b_iN\pmod{p_0}$. The integers $c_i$, $b_i$ and $N$ are prime to $p_0$, by \cref{lem:prime-spec-basic}\ref{item:psb-nonzero} and \ref{item:psb-gcd}. Taking Legendre symbols, we have $\leg{c_i}{p_0}\leg{p_i}{p_0}=\leg{b_i}{p_0}\eps$, and multiplying by $\leg{c_i}{p_0}=\leg{c_i}{p_0}^{-1}$, we obtain \ref{item:rec-pivot}. Part \ref{item:rec-reverse} is quadratic reciprocity for the distinct odd primes $p_0$ and $p_i$.

For \ref{item:rec-cross}, we have $a_iM+b_iN=c_ip_i\equiv0\pmod{p_i}$, that is, $b_iN\equiv-a_iM\pmod{p_i}$. Then
\begin{equation}
b_iL_j=a_jb_iM+b_j(b_iN)\equiv(a_jb_i-a_ib_j)M=\Delta_{ji}M\pmod{p_i}.
\end{equation}
Inserting $L_j=c_jp_j$ and $M=c_0p_0$, we obtain $b_ic_jp_j\equiv\Delta_{ji}c_0p_0\pmod{p_i}$. All factors are prime to $p_i$ by \cref{lem:prime-spec-basic}\ref{item:psb-nonzero}, and taking Legendre symbols and multiplying by $\leg{b_ic_j}{p_i}$, we find \eqref{eq:rec-cross}.
\end{proof}

The characters of fixed integers are determined by residues, by the following consequence of \eqref{eq:reciprocity}.

\begin{lemma}[Characters of fixed integers]\label{lem:fixed-characters}
Let $z=(-1)^\sigma2^e\prod_qq^{f_q}$ be a nonzero integer, where $\sigma\in\{0,1\}$, $e\ge0$, and $q$ runs over the odd primes dividing $z$. Let $p$ be an odd prime not dividing $z$. Then
\begin{equation}\label{eq:fixed-characters}
\chi_p(z)=\sigma\chi_p(-1)+e\chi_p(2)+\sum_qf_q\bigl(\chi_q(p)+\chi_p(-1)\chi_q(-1)\bigr).
\end{equation}
In particular, $\chi_p(z)$ depends only on the residues of $p$ modulo $8$ and modulo the odd primes that divide $z$ to an odd power.
\end{lemma}

\begin{proof}
The map $\chi_p$ is additive, and \eqref{eq:reciprocity} gives $\chi_p(q)=\chi_q(p)+\chi_p(-1)\chi_q(-1)$ for each odd prime $q\ne p$. This proves \eqref{eq:fixed-characters}. By \eqref{eq:reciprocity}, $\chi_p(-1)$ and $\chi_p(2)$ depend only on $p$ modulo $8$, and $\chi_q(p)$ depends only on $p$ modulo $q$. The terms with $f_q$ even vanish in $\F_2$.
\end{proof}

At a prime specialization, we define the following bits. For $1\le i\le11$, let
\begin{equation}\label{eq:theta}
\theta_i=\lambda_0(b_ic_i)+\lambda_0(-1)\lambda_i(-1).
\end{equation}
For distinct $i,j\in\{0,\dots,11\}$, let
\begin{equation}\label{eq:gamma}
\gamma_{ij}=\begin{cases}
\lambda_0(b_jc_j)&\text{if }i=0,\\
\theta_i&\text{if }j=0,\\
\lambda_i(c_0c_jb_i\Delta_{ji})+\theta_i&\text{if }i,j\ge1,
\end{cases}
\end{equation}
and for $\star\in\{E,E'\}$ and $i\in I_F$, let
\begin{equation}\label{eq:gamma-diagonal}
\gamma^\star_{ii}=\lambda_i(\kappa_\star)+\sum_{j\in G_\star\setminus\{i\}}\gamma_{ij}.
\end{equation}

\begin{corollary}[The moving rows as functions of $\eps$]\label{cor:eta-structure}
Let $(M,N)$ be a prime specialization, and let $\eps=(-1)^\eta$ be as in \eqref{eq:eps-def}. Then $\lambda_i(p_j)=\gamma_{ij}+\eta$ for all distinct $i,j\in\{0,\dots,11\}$. For $\star\in\{E,E'\}$ and $i\in I_F\cup O_\star$, the moving row $r^\star_i$ has entry $\lambda_i(g)$ at each constant column $g$, entry $\gamma_{ij}+\eta$ at each moving column $j\ne i$, and, if $i\in I_F$, entry $\gamma^\star_{ii}+\eta$ at the moving column $i$. In other words,
\begin{equation}\label{eq:eta-structure}
r^\star_i=r^\star_i(0)+\eta u_\star,
\end{equation}
where $r^\star_i(0)\in\F_2^{n_\star}$ is the vector with entries $\lambda_i(g)$ at the constant columns, $\gamma_{ij}$ at the moving columns $j\ne i$, and $\gamma^\star_{ii}$ at the moving column $i$ if $i\in I_F$.
\end{corollary}

\begin{proof}
Let $j\ge1$. Then \cref{lem:reciprocity}\ref{item:rec-pivot} implies that $\lambda_0(p_j)=\lambda_0(b_jc_j)+\eta=\gamma_{0j}+\eta$. Let $i\ge1$. By \cref{lem:reciprocity}\ref{item:rec-reverse} and \eqref{eq:reciprocity}, we have $\lambda_i(p_0)=\lambda_0(p_i)+\lambda_0(-1)\lambda_i(-1)$, which is $\theta_i+\eta=\gamma_{i0}+\eta$ by the case just treated. For distinct $i,j\ge1$, we find from \eqref{eq:rec-cross} that
\begin{equation}
\lambda_i(p_j)=\lambda_i(c_0c_jb_i\Delta_{ji})+\lambda_i(p_0)=\lambda_i(c_0c_jb_i\Delta_{ji})+\theta_i+\eta=\gamma_{ij}+\eta.
\end{equation}
This proves the first statement. The entries of $r^\star_i$ at the moving columns $j\ne i$ are $\lambda_i(p_j)=\gamma_{ij}+\eta$, by \cref{def:moving-rows}. For $i\in I_F$, the entry at the moving column $i$ is
\begin{equation}
\lambda_i(\kappa_\star)+\sum_{j\in G_\star\setminus\{i\}}(\gamma_{ij}+\eta)=\gamma^\star_{ii}+7\eta=\gamma^\star_{ii}+\eta,
\end{equation}
since $G_\star$ has eight elements and contains $i$. The entries at the constant columns do not involve $\eta$. Since $u_\star$ has entry $1$ exactly at the moving columns, this proves \eqref{eq:eta-structure}.
\end{proof}

\begin{remark}[Which characters must be controlled]\label{rem:constant-characters}
The bits $\lambda_i(g)$, $\lambda_i(\kappa_\star)$, $\theta_i$, $\gamma_{ij}$ and $\gamma^\star_{ii}$ are characters at $p_0$ or $p_i$ of the fixed integers $-1$, $g\in\Gamma_\star$, $\kappa_\star$, $b_jc_j$ and $c_0c_jb_i\Delta_{ji}$, together with the products $\lambda_0(-1)\lambda_i(-1)$. By \cref{lem:fixed-characters}, they are determined by the residues of $p_0,\dots,p_{11}$ modulo $8$ and modulo the odd primes that divide one of these integers to an odd power; in \cref{sec:rank5,sec:rank6,sec:rank7}, a congruence condition on $(M,N)$ fixes the residues that are needed. Only the pairs $(i,j)$ with $i\in I_F\cup O_\star$ and $j\in G_\star\setminus\{i\}$ occur in the moving rows of $\star$. For $\star=E$, these pairs lie in $I_F\times I_F$, $I_F\times I_H$, $I_J\times I_F$ or $I_J\times I_H$, and for $\star=E'$, they lie in $I_F\times I_F$, $I_F\times I_J$, $I_H\times I_F$ or $I_H\times I_J$. No pair with both indices in $I_H$, or both in $I_J$, occurs. The determinants $\Delta_{ji}$ of such pairs do not enter any moving row, and an odd prime that divides the integers above, for the pairs that occur, only to even powers needs no congruence condition.
\end{remark}

\begin{remark}[The three families]\label{rem:three-families}
The results of this section make no assumption on $p_0$ modulo $4$ or on the form $L_1$, and they apply to the three families without change. In the families of \cref{sec:rank5,sec:rank6,sec:rank7}, we have $c_0=54$, $1$ and $18$, and $p_0\equiv1$, $3$ and $1\pmod8$, respectively. Then $\lambda_0(-1)=0$ and $\theta_i=\lambda_0(b_ic_i)$ in the rank-five and rank-seven families, while $\theta_i=\lambda_0(b_ic_i)+\lambda_i(-1)$ in the rank-six family. In the rank-six and rank-seven families, $L_1=N$ and $c_1=1$, so that $p_1=N$ and $\eps=\leg{p_1}{p_0}$.
\end{remark}

\section{The rank criterion}\label{sec:families-criterion}

With the vectors $r_i^\star(0)$ held fixed, \cref{cor:eta-structure} shows that changing $\eps$ adds the same vector $u_\star$ to every moving row. The following lemma gives a condition under which this does not change the kernel. Its hypothesis cannot be omitted: a moving row equal to $u_\star$ becomes zero when $\eta$ changes. A related common-kernel calculation for two quadratic-residue Laplacians is given in \cite[Lemma~4.2]{Savoie2026}.

\begin{lemma}[Both signs give the same kernel]\label{lem:eta-rank}
Let $\mathbf M(0)$ and $\mathbf M(1)$ be matrices over $\F_2$ with $n$ columns, whose rows are indexed by the same set $R=R_{\mathrm f}\sqcup R_{\mathrm m}$ of fixed and moving rows. Write $r_k(\eta)$ for the row of $\mathbf M(\eta)$ with index $k$. Suppose that there exists $u\in\F_2^n$ such that $r_k(1)=r_k(0)$ for $k\in R_{\mathrm f}$ and $r_k(1)=r_k(0)+u$ for $k\in R_{\mathrm m}$, and that
\begin{equation}\label{eq:u-representation}
u=\sum_{k\in K}r_k(0)
\end{equation}
for a set $K\subseteq R$ such that $K\cap R_{\mathrm m}$ has an even number of elements. Then $\mathbf M(0)$ and $\mathbf M(1)$ have the same row space. In particular, $\ker\mathbf M(1)=\ker\mathbf M(0)$ and $\rk\mathbf M(1)=\rk\mathbf M(0)$.
\end{lemma}

\begin{proof}
By the hypotheses, we have
\begin{equation}
\sum_{k\in K}r_k(1)=\sum_{k\in K}r_k(0)+|K\cap R_{\mathrm m}|\,u=u,
\end{equation}
since $|K\cap R_{\mathrm m}|$ is even. Together with \eqref{eq:u-representation}, this shows that $u$ lies in the row spaces of both matrices. Each row of $\mathbf M(1)$ is the corresponding row of $\mathbf M(0)$, or that row plus $u$, and it lies in the row space of $\mathbf M(0)$. By symmetry, the two row spaces are equal. The kernel of a matrix consists of the vectors $\mathbf w$ with $r\cdot\mathbf w=0$ for every $r$ in its row space. Hence the kernels are equal, and so are the ranks.
\end{proof}

We now define the matrices of the rank criterion. Let $\cP$ be a set of prime specializations, and let $\star\in\{E,E'\}$. Suppose that we are given vectors $f^\star_1,\dots,f^\star_{m_\star}\in\F_2^{n_\star}$, and that the bits
\begin{equation}\label{eq:constant-bits}
\lambda_i(g)\ (i\in I_F\cup O_\star,\ g\in\Gamma_\star),\qquad\lambda_i(\kappa_\star)\ (i\in I_F),\qquad\gamma_{ij}\ (i\in I_F\cup O_\star,\ j\in G_\star\setminus\{i\})
\end{equation}
take the same values at every $(M,N)\in\cP$. Then the vectors $r^\star_i(0)$ of \cref{cor:eta-structure} are the same at every $(M,N)\in\cP$. For $\eta\in\F_2$, we let $\mathbf M_\star(\eta)$ be the matrix with $n_\star$ columns whose rows are
\begin{equation}\label{eq:matrix-def}
f^\star_1,\ \dots,\ f^\star_{m_\star},\ \text{then }r^\star_i(0)+\eta u_\star\ \text{for }i\in I_F,\ \text{then }r^\star_i(0)+\eta u_\star\ \text{for }i\in O_\star,
\end{equation}
with $i$ increasing in each group. We call the first $m_\star$ rows the fixed rows and the last eight rows the moving rows of $\mathbf M_\star(\eta)$. By \cref{cor:eta-structure}, at a point of $\cP$ with $\eps=(-1)^\eta$, the moving rows of $\mathbf M_\star(\eta)$ are the moving rows $r^\star_i$ of \cref{def:moving-rows}.

\begin{proposition}[Rank criterion]\label{prop:rank-criterion}
Let $\datum$ be a twelve-form datum, let $S$ be a set of primes as in \cref{def:prime-specialization}, and let $\cP$ be a set of prime specializations. Suppose that, for $\star=E$ and for $\star=E'$, the list $\Gamma_\star$ consists of $-1$ and distinct primes, including every prime dividing $\kappa_\star$, and that the following hold.
\begin{enumerate}[label=(\roman*),ref=(\roman*)]
\item\label{item:crit-fixed} \emph{Fixed rows.} There exist vectors $f^\star_1,\dots,f^\star_{m_\star}\in\F_2^{n_\star}$ such that, for every $(M,N)\in\cP$, each $f^\star_k$ is a row for $\star_{M,N}$ and $\cB_\star$ in the sense of \cref{def:rows}.
\item\label{item:crit-characters} \emph{Constant characters.} The bits \eqref{eq:constant-bits} take the same values at every $(M,N)\in\cP$. Let $\mathbf M_\star(0)$ and $\mathbf M_\star(1)$ be the matrices \eqref{eq:matrix-def}.
\item\label{item:crit-points} \emph{Points.} There exists an integer $k_\star\ge0$ such that, for every $(M,N)\in\cP$, the image $\delta_\star(\star_{M,N}(\Q))$ contains $k_\star$ linearly independent classes.
\item\label{item:crit-rank} \emph{Rank.} We have $\rk\mathbf M_\star(0)=n_\star-k_\star$, and $u_\star$ is the sum of a set of rows of $\mathbf M_\star(0)$ that contains an even number of moving rows.
\end{enumerate}
Then every $(M,N)\in\cP$ satisfies $d_E=k_E$, $d_{E'}=k_{E'}$ and
\begin{equation}\label{eq:rank-criterion}
\rk E_{M,N}(\Q)=k_E+k_{E'}-2.
\end{equation}
Moreover, $\delta_\star(\star_{M,N}(\Q))$ consists of the classes $\sq{\mathbf w}_{\cB_\star}$ with $\mathbf w\in\ker\mathbf M_\star(0)$.
\end{proposition}

\begin{proof}
Fix $(M,N)\in\cP$ and $\star\in\{E,E'\}$, and let $\eta\in\F_2$ be such that $\eps=(-1)^\eta$ at $(M,N)$. By \cref{lem:prime-spec-basic}, the list $\cB_\star$ is an ambient basis for $\star$. We first show that every row of $\mathbf M_\star(\eta)$ is a row in the sense of \cref{def:rows} at $(M,N)$. For the fixed rows, this is hypothesis \ref{item:crit-fixed}. By hypothesis \ref{item:crit-characters} and \cref{cor:eta-structure}, the moving rows of $\mathbf M_\star(\eta)$ are the rows $r^\star_i$ of \cref{def:moving-rows} at $(M,N)$, and these are rows by \cref{lem:moving-rows}. By \cref{lem:upper-bound}, the coordinate vectors of the classes in $\delta_\star(\star(\Q))$ lie in $\ker\mathbf M_\star(\eta)$.

The matrices $\mathbf M_\star(0)$ and $\mathbf M_\star(1)$ satisfy the hypotheses of \cref{lem:eta-rank}, with $u=u_\star$, by \eqref{eq:matrix-def} and hypothesis \ref{item:crit-rank}. This lemma implies that $\ker\mathbf M_\star(\eta)=\ker\mathbf M_\star(0)$ and $\rk\mathbf M_\star(\eta)=n_\star-k_\star$. By hypothesis \ref{item:crit-points}, there are $k_\star$ linearly independent classes in $\delta_\star(\star(\Q))$, and their coordinate vectors are linearly independent. Applying \cref{cor:exact-dimension} to $\mathbf M_\star(\eta)$, we find that $d_\star=k_\star$ and that $\delta_\star(\star(\Q))$ consists of the classes $\sq{\mathbf w}_{\cB_\star}$ with $\mathbf w\in\ker\mathbf M_\star(\eta)=\ker\mathbf M_\star(0)$. Finally, by \cref{lem:rank-formula}, we have $\rk E_{M,N}(\Q)=d_E+d_{E'}-2=k_E+k_{E'}-2$.
\end{proof}

In \cref{sec:rank5,sec:rank6,sec:rank7}, hypothesis \ref{item:crit-fixed} is a lemma on the local images at the fixed places, hypothesis \ref{item:crit-characters} follows from a congruence condition on $(M,N)$ and \cref{rem:constant-characters}, hypothesis \ref{item:crit-points} is verified with explicit points and \cref{lem:independence}, and hypothesis \ref{item:crit-rank} is a finite computation with a certificate that can be checked by hand. Since the criterion holds for both values of $\eps$, the prime specializations produced in \cref{sec:primes} need not satisfy any condition on $\leg N{p_0}$.
